% SPDX-License-Identifier: CC-BY-4.0
% Copyright (c) 2026 Martin Schröder <info@swedishembedded.com>

\section{Related work}\label{sec:related}

\emph{Self-training on self-proposed problems.} The closest prior design is Haluptzok et al.: a model
synthesises programming problems with solutions, an interpreter filters them for correctness, and the
model is fine-tuned on the result.~\citep{haluptzok2022}\cl{CT24} STaR and ReST-EM also restart from the
base each iteration on accumulated, filtered self-generated data, but they filter the model's
\emph{outputs} against known answers;~\citep{star2022,restem2023}\cl{CT02,CT16} here the model generates
the \emph{inputs} and the oracle supplies the outputs. Self-Instruct generates inputs too, but labels them
with the same model;~\citep{selfinstruct2022}\cl{CT01} ReST iterates grow and improve steps on filtered
samples.~\citep{rest2023} Noisy Student iterates generations on a growing pool, with a teacher's
labels.~\citep{noisystudent2019}\cl{CT26} Expert iteration amplifies the policy by search and imitates
the result;~\citep{exit2017}\cl{CT19} an interpreter is ground truth, not an amplified policy, so only the
generate--filter--retrain skeleton is shared. In self-play designs a proposer is trained to set
learnable tasks, as in asymmetric self-play and in Absolute Zero, where a code executor validates the
proposed tasks;~\citep{asymselfplay2017,absolutezero2025}\cl{CT28,CT18} our proposer is not trained to
propose, only filtered. Other self-improvement work changes a model's own context or scaffold rather than
its training data: STOP improves the improver scaffold itself,~\citep{stop2023}\cl{CT25} Learning to
Self-Evolve trains a model to improve its context at test time,~\citep{lse2026}\cl{CT31} and hyperagents
let a meta agent modify the agent;~\citep{hyperagents2026}\cl{CT32} SEAL has a model write its own
fine-tuning data and directives.~\citep{seal2025}\cl{CT33} Our matched-data control does not show that the
recursion itself helps here. On reinforcement learning with verifiable rewards, Yue et al.\ find gains at
small $k$ bounded by the base model at large $k$, and new patterns from distillation;~\citep{rlvrbound2025}\cl{CT34}
our own RL arms are too narrow to add to that question.

\emph{Selection and curricula.} Selecting what to label is active learning, here by query synthesis with
the model's own disagreement as the acquisition function; ordering by difficulty is curriculum
learning;~\citep{settles2012,curriculum2009}\cl{CT29} learned curricula over problem categories are a
recent variant.~\citep{sec2025}\cl{CT23}

\emph{Predicting execution.} Predicting what a program prints is CRUXEval's output-prediction
task;~\citep{cruxeval2024}\cl{CT17} models have been trained to predict execution since Learning to
Execute, and more recently with execution traces.~\citep{learningtoexecute2014,codeexecutor2023,semcoder2024}\cl{CT27}
CodeRL trains a generator with a critic of functional correctness from unit tests;~\citep{coderl2022}\cl{CT20}
training repeatedly on a model's own generated content can make the tails of the original distribution
disappear,~\citep{collapse2023}\cl{CT21} which our series counter only by labelling every program with the
interpreter.

\emph{Evaluation.} A greedy gate that accepts any higher validation score lets inferior updates through
across rounds; TRUST controls their rate by online multiple testing.~\citep{trust2026}\cl{CT15} Reusing a
holdout adaptively can overfit it,~\citep{holdout2015}\cl{CT22} and fresh test sets test
replication;~\citep{recht2019} evaluations need clustered standard errors and seed-level
variance.~\citep{errorbars2024,benchvariance2021,precipice2021}\cl{CT30} Related work for the system
components is in \cref{app:related-system}.
